\chapter{Réalisation du backend}
\label{ch:backend}

\section{Architecture backend}

Le serveur suit un découpage en quatre couches appliqué sans exception à l'ensemble des
modules : un contrôleur REST expose les endpoints et porte les règles d'autorisation ; un
service concentre la logique métier et la démarcation transactionnelle ; un repository
Spring~Data assure la persistance ; des DTO, implémentés sous forme de \code{record} Java,
transportent les données à l'entrée et à la sortie sans jamais exposer les entités JPA.

Cette dernière règle est importante : aucune entité n'est sérialisée directement en JSON.
Elle évite deux écueils classiques --- la divulgation involontaire d'attributs sensibles et
les initialisations paresseuses déclenchées pendant la sérialisation --- au prix d'une
étape de conversion, assurée par des mappers écrits à la main (\code{MemberMapper}, \code{GovernanceMapper}).

Le volume réalisé s'établit à 22 contrôleurs, 28 services, 30 repositories et 30 entités,
soit environ 10~000 lignes de Java.

\section{Organisation des packages}

Le découpage des packages est fonctionnel et non technique : on ne trouve pas un package
\code{controllers} et un package \code{services}, mais un package par domaine métier
contenant ses propres couches. Le tableau~\ref{tab:packages} en donne la structure.

\begin{table}[H]
\centering
\caption{Organisation des packages métier du backend}
\label{tab:packages}
\small
\begin{tabularx}{\textwidth}{@{}p{3.2cm}X@{}}
\toprule
\textbf{Package} & \textbf{Responsabilité} \\
\midrule
\code{auth} & Comptes, authentification, jetons de réinitialisation, historique des rôles \\
\code{members} & Profils, statuts, périodes d'affiliation \\
\code{governance} & Équipes, mandats, appartenance aux équipes, dérivation des rôles \\
\code{membershiprequest} & Candidatures d'adhésion et décisions \\
\code{event} & Événements, organisateurs, participants \\
\code{publication} & Publications et co-signataires \\
\code{meeting} & Réunions et procès-verbaux \\
\code{document} & Archivage documentaire et versions \\
\code{equipment} & Inventaire, distributions, demandes de matériel \\
\code{partnership} & Partenaires et conventions \\
\code{notification} & Notifications, discussions, messages \\
\code{calendar} & Agrégation du calendrier \\
\code{report} & Rapports annuels et génération PDF \\
\code{audit} & Journal d'audit \\
\code{administration} & Administration des comptes, recherche globale \\
\code{laboratory} & Identité institutionnelle \\
\code{security} & Filtre JWT, détails utilisateur, service de jetons \\
\code{common} & Classe de base des entités, stockage de fichiers, gestion des erreurs, énumérations \\
\bottomrule
\end{tabularx}
\end{table}

Ce découpage présente un avantage pratique vérifié pendant le projet : ajouter un module
n'oblige à toucher aucun package existant, hormis l'enregistrement éventuel d'une règle
d'autorisation partagée.

\section{Authentification}

L'authentification est sans état. La méthode \code{login} vérifie les identifiants via le
gestionnaire d'authentification de Spring~Security, puis émet deux jetons : un jeton
d'accès valable une heure et un jeton de rafraîchissement valable sept jours. Les mots de
passe sont hachés avec BCrypt au facteur de coût~12.

Le filtre \code{JwtAuthenticationFilter} s'exécute avant la chaîne d'autorisation. Il
extrait le jeton, en vérifie la signature et le type, puis --- et c'est le point
déterminant --- recharge l'utilisateur depuis la base par
\code{CustomUserDetailsService.loadUserByUsername}. Les habilitations attachées au contexte
de sécurité proviennent donc toujours de l'état courant en base, jamais de la charge utile
du jeton. Le filtre vérifie en outre que le compte est actif, ce qui permet à un gel de
statut de prendre effet immédiatement.

Le module comprend également la réinitialisation de mot de passe par courriel, adossée à
l'entité \code{PasswordResetToken} (jeton à usage unique avec date d'expiration), et le
changement de mot de passe par l'utilisateur connecté.

\section{Gestion des utilisateurs et des membres}

La distinction entre \code{User} et \code{Member} structure ce module. \code{MemberService}
porte les opérations métier : consultation paginée avec filtres, mise à jour de profil,
changement de statut, départ, retour, archivage et restauration.

La règle d'édition du profil mérite d'être citée car elle illustre une autorisation qui ne
se réduit pas à un rôle : un membre peut modifier son propre profil, et seuls
l'administrateur, le directeur et le vice-directeur peuvent modifier celui d'un autre. Ce
contrôle est porté par le service et non par une annotation, parce qu'il dépend de
l'identité de la cible et non seulement du rôle de l'appelant.

L'archivage est une suppression logique. La classe \code{BaseEntity} porte les champs
\code{deleted} et \code{deletedAt}, et l'annotation Hibernate
\code{@SQLRestriction("deleted = false")} exclut automatiquement les lignes archivées de
toutes les requêtes. Cette élégance a une contrepartie : une entité archivée encore
référencée par une autre devient inaccessible et lève une exception lors de
l'initialisation paresseuse. L'utilitaire \code{SoftDeleteSafe} encapsule ces accès et
substitue une valeur de repli --- « Membre archivé », « Matériel archivé » --- plutôt que
de laisser l'erreur remonter et invalider toute la réponse.

\section{Séparation \srt{} / \rol{} / \man{}}
\label{sec:srm}

Ce service constitue le cœur du projet. \code{UserRoleSyncService} applique un unique
principe : le rôle applicatif n'est pas une donnée saisie mais une fonction du statut et
des mandats actifs. Le listing~\ref{lst:compute} en donne la règle de résolution.

\begin{lstlisting}[language=Java,caption={Résolution du rôle effectif (\code{UserRoleSyncService})},label={lst:compute}]
private RoleName computeEffectiveRole(Member member) {
    MemberStatus status = member.getCurrentStatus();
    if (status == MemberStatus.RETIRED || status == MemberStatus.FORMER) {
        // Un compte gele ne conserve jamais de permissions de gouvernance,
        // meme si un mandat n'a pas ete explicitement clos.
        return STATUS_TO_ROLE.get(status);
    }
    List<Mandate> activeMandates =
        mandateRepository.findByMemberIdAndActiveTrueAndEndDateIsNull(member.getId());
    for (MandateRoleType type : MANDATE_PRIORITY) {   // DIRECTOR > VICE_DIRECTOR > TEAM_LEADER
        if (activeMandates.stream().anyMatch(m -> m.getRole() == type)) {
            return MANDATE_TO_ROLE.get(type);
        }
    }
    return STATUS_TO_ROLE.getOrDefault(status, RoleName.ASSOCIATE_MEMBER);
}
\end{lstlisting}

Trois propriétés de cette fonction sont vérifiées par les tests du
chapitre~\ref{ch:tests}. Le statut gelé l'emporte sur tout, y compris sur un mandat qui
n'aurait pas été clos --- c'est une défense en profondeur, puisque le changement de statut
clôt déjà les mandats. La priorité entre mandats est stricte et totale. Enfin, en l'absence
de mandat, le rôle retombe sur celui que le statut implique.

Toute application effective d'un rôle passe par \code{applyRole}, qui clôt la ligne
d'historique ouverte, en ouvre une nouvelle avec sa cause --- changement de mandat,
changement de statut, dérogation administrative ou réconciliation --- puis écrit le rôle
sur le compte. Si le rôle calculé est déjà celui enregistré et déjà historisé, la méthode
ne fait rien : l'opération est idempotente.

Ce service n'est cependant invoqué que sur les chemins d'écriture. Un compte créé avant
son introduction, et jamais modifié depuis, conserverait un rôle obsolète. C'est
exactement ce qui a été constaté lors de l'audit fonctionnel. La réponse a été
\code{reconcileAll()}, exécuté au démarrage par un \code{CommandLineRunner} : la méthode
parcourt les membres, recalcule le rôle attendu, corrige les écarts et complète les
historiques manquants. Sur un système déjà cohérent, elle n'écrit rien et le journal
indique simplement que tous les comptes sont conformes.

\section{Gouvernance}

\code{MandateService} implémente la règle d'unicité du titulaire énoncée au cahier des
charges : attribuer un mandat de directeur ou de vice-directeur clôt automatiquement celui
du titulaire précédent. Le mandat n'est jamais supprimé --- sa date de fin est renseignée
et son indicateur d'activité abaissé --- de sorte que l'historique de gouvernance reste
intégralement reconstituable. Chaque attribution et chaque clôture déclenche la
resynchronisation des rôles décrite en~§\ref{sec:srm} et illustrée par la
figure~\ref{fig:seqMandat}.

L'attribution et la clôture d'un mandat sont réservées à l'administrateur système et au
directeur. Le vice-directeur, bien qu'il accède au module de gouvernance, n'y dispose que
d'une consultation et du transfert d'un membre entre équipes.

\section{Équipes}

\code{TeamService} gère les équipes de recherche, dont le responsable n'est pas un champ de
la table mais une donnée déduite du mandat \code{TEAM\_LEADER} actif rattaché à l'équipe.
Ce choix évite une incohérence possible entre un champ « chef d'équipe » et la table des
mandats.

L'appartenance d'un membre à une équipe est elle aussi historisée :
\code{TeamMembershipHistory} conserve chaque période, et le transfert d'un membre clôt la
période courante avant d'en ouvrir une nouvelle.

\section{Équipements}

Le module matériel répond aux sections~14 et~15 du cahier des charges et repose sur un
invariant simple : la quantité disponible d'un article diminue exactement du nombre
d'unités distribuées et remonte du même nombre au retour. \code{EquipmentService.distribute}
refuse toute distribution excédant le stock disponible, et \code{returnDistribution} refuse
un retour déjà enregistré.

L'exigence la plus spécifique du cahier des charges est le suivi de « qui n'a pas reçu ».
Elle est satisfaite par une requête JPQL dédiée, \code{findMembersWithoutDistribution},
qui calcule le complément : les membres actifs ne figurant pas parmi les bénéficiaires
d'un article donné. Les demandes de matériel suivent un cycle propre --- soumission par
tout membre des modules internes, validation ou rejet par le directeur ou le
vice-directeur --- avec notification des décideurs à la soumission.

\section{Documents et versions}

\code{DocumentService} implémente l'archivage documentaire. Un document porte des
métadonnées --- titre, catégorie parmi neuf valeurs, étiquettes, événement de rattachement
éventuel --- et au moins une version. Le dépôt d'un nouveau fichier n'écrase jamais le
précédent : il ajoute une ligne \code{DocumentVersion} avec son numéro d'ordre, son
déposant et son horodatage.

Le stockage physique est délégué à \code{FileStorageService}, qui range les fichiers selon
une arborescence \code{catégorie/aaaa/MM/jj/} et attribue à chacun un nom aléatoire tout en
conservant le nom d'origine en base. Le répertoire racine est lu depuis la variable
d'environnement \code{FILE\_STORAGE\_DIR}, ce qui rend le déploiement indépendant de la
machine. Le téléchargement passe par un endpoint authentifié qui renvoie le flux avec son
type et son nom d'origine, plutôt que par un lien direct vers le disque.

\section{Événements}

\code{EventService} couvre les conférences, séminaires et workshops. Deux mécanismes
répondent à des exigences précises du cahier des charges : l'attribut \code{editionNumber}
et l'endpoint \code{/events/editions} permettent de retrouver les occurrences successives
d'une même manifestation ; les entités \code{EventOrganizer} et \code{EventParticipant}
distinguent l'équipe organisatrice des inscrits. La création et la modification sont
ouvertes aux chefs d'équipe, la suppression reste réservée à la direction, et l'inscription
est accessible à tout membre authentifié.

\section{Notifications}

\code{NotificationService} expose quatre modes d'émission, correspondant à quatre
politiques de destinataires : notification individuelle, diffusion à une liste,
notification de la gouvernance et diffusion à tous les membres actifs. Le troisième mérite
attention : les destinataires ne sont pas déterminés par un rôle stocké mais par les
mandats de directeur et de vice-directeur \emph{actuellement actifs}. Une notification de
décision suit donc automatiquement l'évolution de la gouvernance, sans configuration.

La diffusion générale exclut les membres retraités et anciens, cohérente avec le fait que
ces comptes n'ont plus d'accès actif.

\section{Audit et traçabilité}

\code{AuditService} enregistre l'acteur, l'action, le type et l'identifiant de l'entité
concernée, un libellé explicatif, l'adresse IP et l'horodatage. Son point notable est
l'usage de la propagation \code{REQUIRES\_NEW} : l'écriture de la trace s'effectue dans une
transaction distincte de l'opération métier. Si celle-ci échoue et est annulée, la trace de
la tentative subsiste --- ce qui est précisément l'objet d'un journal d'audit.

Comme indiqué au chapitre~\ref{ch:conception}, l'entité d'audit ne possède pas
d'association JPA vers les entités tracées. Ce découplage assumé permet d'auditer n'importe
quel type d'objet sans multiplier les clés étrangères ni rendre la trace dépendante de la
survie de sa cible.

\section{Rapports annuels et génération PDF}

\code{ReportService} agrège les indicateurs d'une année --- événements, publications,
activités --- et les conserve dans l'entité \code{AnnualReport}. La production du document
est confiée à \code{AnnualReportPdfGenerator}, bâti sur OpenPDF, qui compose une mise en
page structurée archivée via le service de stockage ; le fichier est ensuite téléchargeable
par un endpoint authentifié.

Le cahier des charges demande une génération « annuelle automatique ». Elle est assurée par
\code{AnnualReportScheduler}, déclenché par une expression \code{cron} configurable dont la
valeur par défaut porte l'exécution au 1\textsuperscript{er} janvier à 02~h, heure de
Casablanca, afin de clore l'année civile écoulée. La génération manuelle depuis l'interface
reste possible et produit le même document, ce qui permet de préparer le rapport avant
l'échéance.
